\section{The Lord Who Entered Death}

Christ's victory begins with a death that is fully real. The eternal Son assumed a complete human nature, with a human soul and body. On Calvary it is the divine Son who dies, and the death he undergoes is human: his soul is separated from his body. His divinity remains immortal. During that separation both soul and body remain personally united to the Word (CCC 624--626). The Incarnation is not undone between Good Friday and Easter. The Lord who lies in the tomb is the same Lord who has spoken to the disciples and who will stand among them risen.

Aquinas explains the continuing union by attending to the person of Christ (ST III, q.~50, aa.~2--3). Death divides his human soul from his human body; it does not divide the one person of the Word into two, nor remove either soul or body from that person. The Creed therefore attributes both burial and descent among the dead to the Son himself. We can trust that the Son has truly entered the condition of the dead and remained united to the humanity he assumed for our salvation. Neither his body in the tomb nor his soul among the dead has been abandoned by the Word.

His death is freely accepted. Jesus speaks of laying down his life and taking it up again (Jn 10:17--18). That freedom does not make the sufferings unreal. In Gethsemane he is sorrowful and distressed; he asks that the chalice pass, and freely submits to the Father's will (Mk 14:32--42). The Letter to the Hebrews remembers his cries and tears, and teaches that the Son learned obedience through what he suffered (Heb 5:7--9). He who is without sin knows the human shrinking before suffering and death, and his human obedience remains entire within it.

We can therefore distinguish fear from refusal. An involuntary recoil before dying does not by itself mean that the will has abandoned God. The Lord's own agony prevents us from requiring ourselves or another person to feel no anguish as a condition of faithful surrender. Faithfulness may take the form of returning, within anguish, to the Father's will. The desire for relief, the request for companionship, and the effort to remain obedient can coexist. Christ himself asks the disciples to remain and watch.

On the cross the Lord continues to care for others. He receives the penitent thief and gives his mother and the beloved disciple to one another (Lk 23:39--43; Jn 19:25--30). While his enemies mock him, he promises paradise to the man beside him and provides for those who remain. His final surrender to the Father completes this obedient love (Lk 23:44--46). The dying Christian can turn to a Savior who has known bodily helplessness and continued to love within it.

When Christ rises, the body that bore the wounds is glorified beyond death's reach. He shows his hands and his side to the disciples (Jn 20:19--29). Their Lord's identity continues through death and resurrection, while his manner of bodily life is transformed. He has done more than return to the mortal condition from which Lazarus was recalled. Paul names him the firstfruits of those who have fallen asleep and joins his Resurrection to the resurrection of those who belong to him (1~Cor 15:20--28).

This is why we can offer our own death in union with his obedience and love (CCC 1011). The offering begins with what Christ has accomplished. We do not make our death redemptive by sufficient intensity of feeling. He draws our frailty into his own saving act. While we are able, we can deliberately entrust the manner, time, and circumstances of dying to the Father, asking for the grace to remain faithful through what we cannot control. Such an intention need not be accompanied by certainty that we will feel brave when the hour arrives.

We entrust our future to this crucified and risen Son. He saves those who die in his grace and brings them into his life. Asking for a good death therefore includes asking for the faith to rely on him when pain or weakness makes our own resources fail.

